\documentclass[a4paper,11pt]{article}
\usepackage{exam-style-341}
\examinehead{341 农业综合知识三 · 程序设计}{模拟试卷（一）参考答案}
\begin{document}
\answercover{341 农业综合知识三 · 程序设计}{模拟试卷（一）参考答案}

\answerpart{一、单项选择题（每题 2 分，共 20 分）}
\begin{qgroup}
\q A。
\begin{analysis}
标识符只能由字母、数字和下划线组成，且首字符不能是数字，也不能使用关键字。\code{\_num2} 合法；\code{2num} 以数字开头；\code{-num} 含减号；\code{double} 是关键字。
\end{analysis}

\q C。
\begin{analysis}
优先级由高到低为：\code{!}（单目逻辑非）$>$ \code{!=}（关系运算符）$>$ \code{\&\&}（逻辑与）$>$ \code{=}（赋值运算符），所以 \code{=} 的优先级最低。
\end{analysis}

\q B。
\begin{analysis}
\code{k = 3} 与 \code{case 3} 匹配，执行 \code{k++} 后 \code{k = 4}；由于该分支后没有 \code{break}，会继续执行 \code{case 4} 中的 \code{k++}，\code{k = 5}，随后遇到 \code{break} 退出 \code{switch}。所以结果为 5。
\end{analysis}

\q B。
\begin{analysis}
\code{continue} 结束本次循环、进入下一次迭代：\code{i} 为 1、3、5 时累加，\code{i} 为 2、4 时被跳过。\code{s = 1 + 3 + 5 = 9}。
\end{analysis}

\q A。
\begin{analysis}
字符串 \code{"China"} 占 5 个字符加上 1 个字符串结束标志 \code{'\textbackslash 0'}，而 \code{s} 是用字符串常量初始化的字符数组，其大小 \code{sizeof(s) = 6}；\code{strlen} 统计的是 \code{'\textbackslash 0'} 之前的字符个数，为 5。
\end{analysis}

\q A。
\begin{analysis}
C 语言函数参数采用值传递，形参 \code{x} 只是实参 \code{a} 的副本，函数内对 \code{x} 的修改不会影响 \code{a}，因此输出 5。
\end{analysis}

\q B。
\begin{analysis}
\code{p} 指向 \code{a[0]}，\code{*(p + 2)} 等价于 \code{a[2]}，值为 5；\code{p[4]} 等价于 \code{a[4]}，值为 9。两者之和为 14。
\end{analysis}

\q A。
\begin{analysis}
共用体的所有成员共用同一段存储单元，变量所占字节数等于最长成员所占的字节数，所以 A 正确。结构体各成员各自占用存储单元，B 错误；共用体变量在定义时只能对第一个成员初始化，C 错误；结构体变量可以作为函数参数（值传递），D 错误。
\end{analysis}

\q A。
\begin{analysis}
\code{"r"} 以只读方式打开一个已存在的文本文件；\code{"w"} 以只写方式打开，若文件存在则清空其内容；\code{"a"} 以追加方式打开，只能写不能读；\code{"rb+"} 以二进制方式打开并允许读写。故只有 A 符合“只读”要求。
\end{analysis}

\q B。
\begin{analysis}
\code{strcat(a, b)} 把 \code{b} 指向的字符串 \code{"World"} 连接到 \code{a} 中字符串 \code{"Hello"} 的末尾，\code{a} 变为 \code{"HelloWorld"}，长度为 10，未越界（\code{a} 大小为 20）。
\end{analysis}
\end{qgroup}

\answerpart{二、填空题（每题 2 分，共 10 分）}
\begin{qgroup}
\q 2；3.5。
\begin{analysis}
\code{a / b} 的两个操作数都是 \code{int}，执行整数除法，结果为 2；\code{a / c} 中 \code{c} 是 \code{double}，\code{a} 先被转换为 \code{double} 再做实数除法，结果为 3.5。
\end{analysis}

\q \code{x \% 7 == 0 \&\& x \% 2 == 0}（也可写成 \code{x \% 14 == 0}）。
\begin{analysis}
“能被 7 整除”即 \code{x \% 7 == 0}，“是偶数”即 \code{x \% 2 == 0}，两个条件必须同时成立，用逻辑与 \code{\&\&} 连接。
\end{analysis}

\q 4；\code{'x'}。
\begin{analysis}
\code{"exam"} 有 4 个可见字符，\code{strlen(str)} 为 4；下标从 0 开始，\code{str[0] = 'e'}，\code{str[1] = 'x'}。
\end{analysis}

\q 2；1。
\begin{analysis}
形参 \code{p}、\code{q} 分别保存主函数中 \code{x}、\code{y} 的地址，函数内通过 \code{*p}、\code{*q} 直接访问并交换了 \code{x}、\code{y} 的值，所以调用后 \code{x = 2}、\code{y = 1}。
\end{analysis}

\q \code{"in.txt"}；\code{fclose}。
\begin{analysis}
\code{fopen} 的第一个参数是文件路径字符串，必须用双引号括起来；文件使用完毕后应用 \code{fclose} 关闭，以把缓冲区中的数据写回文件并释放文件指针。
\end{analysis}
\end{qgroup}

\answerpart{三、程序阅读题（每题 5 分，共 10 分）}
\begin{qgroup}
\q 运行结果为：\code{19 32}
\begin{analysis}
数组 \code{a} 中的奇数为 3、1、1、5、9，共 5 个。函数 \code{fun} 用指针 \code{p} 遍历数组，累加所有奇数，返回 $3+1+1+5+9=19$。主函数中的循环对每一个奇数执行一次 \code{t *= 2}，共执行 5 次，故 $t=2^{5}=32$。两个实参的求值互不影响，\code{printf} 输出 \code{19 32}。
\end{analysis}

\q (1) \code{s[i] != '\textbackslash 0'}（写成 \code{s[i]} 亦可）；\quad (2) \code{s[i] >= '0' \&\& s[i] <= '9'}。
\begin{analysis}
字符串以 \code{'\textbackslash 0'} 作为结束标志，循环必须遍历到该标志之前，故 (1) 填 \code{s[i] != '\textbackslash 0'}。判断数字字符既可写 \code{s[i] >= '0' \&\& s[i] <= '9'}，也可写库函数 \code{isdigit(s[i])}（需包含 \code{<ctype.h>}）。
\end{analysis}
\begin{analysis}
说明：题干程序使用了 \code{gets} 读入整行字符串，该函数已从 C11 标准中删除，实际编程建议改用 \code{fgets(s, 80, stdin)}（此时需另行处理行末的换行符）。
\end{analysis}
\end{qgroup}

\answerpart{四、程序设计题（每题 5 分，共 10 分）}
\begin{qgroup}
\q 参考实现：
\begin{lstlisting}[language={[custom]C}]
#include <stdio.h>
int main(void) {
    int a[10], i, max, min;
    double sum = 0.0;
    printf("请输入 10 个整数：\n");
    for (i = 0; i < 10; i++)
        scanf("%d", &a[i]);
    max = min = a[0];                 /* 以第一个元素作为初值 */
    for (i = 0; i < 10; i++) {
        sum += a[i];
        if (a[i] > max) max = a[i];
        if (a[i] < min) min = a[i];
    }
    printf("max=%d, min=%d, average=%.2f\n", max, min, sum / 10.0);
    return 0;
}
\end{lstlisting}
\begin{analysis}
评分要点：正确定义数组并读入 10 个整数（1 分）；\code{max}、\code{min} 用 \code{a[0]} 初始化并在同一轮循环中更新（2 分）；用 \code{double} 型变量累加求和，除以 \code{10.0} 得到实数平均值（1 分）；用 \code{\%.2f} 按要求输出两位小数（1 分）。
\end{analysis}
\begin{analysis}
易错点：\code{max}、\code{min} 若初始化为 0，当输入全为正数或全为负数时会得到错误结果，应初始化为数组首元素；求平均值时若写成 \code{sum / 10} 且 \code{sum} 为整型，会发生整数除法截断。
\end{analysis}

\q 参考实现：
\begin{lstlisting}[language={[custom]C}]
#include <stdio.h>
int is_prime(int n) {                 /* 是素数返回 1，否则返回 0 */
    int i;
    if (n < 2) return 0;
    for (i = 2; i * i <= n; i++)
        if (n % i == 0) return 0;
    return 1;
}
int main(void) {
    int n, count = 0;
    for (n = 100; n <= 200; n++) {
        if (is_prime(n)) {
            printf("%5d", n);
            count++;
            if (count % 5 == 0) printf("\n");
        }
    }
    if (count % 5 != 0) printf("\n");
    return 0;
}
\end{lstlisting}
\begin{analysis}
评分要点：\code{is\_prime} 函数对 \code{n < 2} 的处理（1 分）；用 \code{i * i <= n}（或 \code{i <= sqrt(n)}）控制试除范围并在整除时提前返回 0（2 分）；主函数遍历 100\textasciitilde 200 并调用函数（1 分）；用计数器 \code{count} 控制每行 5 个输出（1 分）。
\end{analysis}
\begin{analysis}
运行结果：101、103、107、109、113、127、131、137、139、149、151、157、163、167、173、179、181、191、193、197、199，共 21 个素数，前 20 个分 4 行输出，最后一个 \code{199} 单独占第 5 行。若在循环中写 \code{i <= n / 2} 或 \code{i < n} 也不会出错，但效率较低；使用 \code{sqrt} 时需要包含 \code{<math.h>}。
\end{analysis}
\end{qgroup}

\end{document}
